\documentclass[11pt]{article}
\usepackage{icmlko}
\begin{document}

\title{의심이 틀렸고 헤드라인은 더 강해졌다 --- 다만 ``최적점은 다섯''은 안 선다}
\author{월드모델 실험실}
\date{노트 618 · 619 · 2026년 8월 5일}
\maketitle

\begin{abstract}
이 저장소에는 갈래가 둘이다 --- 교차도메인 순위 모형(\texttt{lab/})과
팝업 조건$\to$결과 예측(\texttt{harness/} $\cdot$ \texttt{state/}). 이번
세션이 \texttt{lab/} 쪽에서 만든 자(표본 부트스트랩 $\cdot$ 군집 $\cdot$
유의성 위 선택)를 \textbf{팝업 갈래의 헤드라인에 겨눴다}.

헤드라인은 논문 $4$ 다 --- ``계수 계열 $n{=}75$ 에서 피처를 $132\!\to\!5$ 로
깎자 \textbf{처음으로 상수를 유의하게 이겼다}''($\Delta=-0.0540$,
CI$[-0.103,-0.010]$). 의심이 둘이었다: \textbf{피처 수를 다섯 개 훑어
최댓값을 골랐고}, CI 가 \textbf{행 i.i.d.} 인데 폴드는 IP 무결이다.

\begin{center}
\begin{tabular}{lrrll}
\hline
피처 & $p$ & $\Delta$ & IP 군집 $95\%$ & 군집 $99\%$(본페로니) \\
\hline
$132$ 전체 & $132$ & $-0.033$ & $0$ 을 문다 & $0$ 을 문다 \\
$10$ 기획 & $20$ & $-0.040$ & $\mathbf{0}$ \textbf{밖} & $0$ 을 문다 \\
\textbf{5 기여} & $10$ & $\mathbf{-0.058}$ & $[-0.098,-0.019]$ & $\mathbf{[-0.114,-0.008]}$ \\
\textbf{3 상위} & $6$ & $-0.057$ & $\mathbf{0}$ \textbf{밖} & $\mathbf{[-0.100,-0.015]}$ \\
\hline
\end{tabular}
\end{center}

\textbf{의심 둘 다 빗나갔다.} 군집 구간은 $0$ 밖이고 본페로니 $99\%$ 도
$0$ 밖이다. \textbf{그런데 다른 것이 무너졌다} --- 논문 $4$ 의 ``셋으로
줄이면 다시 못 이겼다, 최적점은 다섯''이 재현되지 않는다.
\end{abstract}

\begin{figure}[t]
\centering
\includegraphics[width=\textwidth]{figs/f1.pdf}
\caption{\textbf{가} --- 다섯도 셋도 본페로니를 통과한다.
\textbf{나} --- 군집이 안 문다(행/IP $=1.05$).
\textbf{다} --- 두 갈래가 같은 다섯 축으로 수렴했다.}
\end{figure}

\section{군집 의심이 왜 빗나갔나}

폴드를 IP 무결로 나누는 판에서 구간만 행 i.i.d. 로 내는 것은 원리상
틀렸다(노트 $608$ 이 \texttt{lab/} 쪽에 ``안 했다''고 적어 둔 항목이다).
그런데 재 보니 구간 폭이 $0.89\!\sim\!0.99$ 배로 \textbf{거의 안 변한다.}

\begin{center}
채점 $40$ 행에 IP $38$ 개 --- \textbf{행/IP $=1.05$}
\end{center}

\textbf{묶을 것이 없으면 군집 보정이 할 일이 없다.} 이 풀은 사실상
IP 하나에 행 하나다. \textbf{조항은 여전히 옳고 이 자료에서만 안 문다} ---
시장 레코드가 붙어 IP 당 행이 늘면 그때는 문다.

\section{선택 의심은 본페로니로 갚았다}

다섯 집합을 훑어 최댓값을 골랐으니 명목 $95\%$ 는 못 쓴다. 각 검정을
$\alpha{=}0.01$ 로 조여도 \textbf{$5$ 기여 속성이 $[-0.114,-0.008]$ 로
$0$ 밖}이다. 논문 $4$ 의 결론은 \textbf{다중검정을 감안하고도 선다.}

\begin{quote}
\textbf{내가 코드로 낸 ``선택 $p{=}0.619$''는 유효한 검정이 아니다.}
관측 통계량을 부트스트랩해 그 최소값 분포와 견줬는데, 그 분포는 관측치
근처에 놓이므로 확률이 $0.5$ 언저리로 나오는 것이 \emph{정의상} 당연하다.
영가설 분포가 아니다. \textbf{다중검정 보정은 본페로니 구간이 맡고 그것은
통과한다.}
\end{quote}

\section{무너진 것 --- ``최적점은 다섯''}

논문 $4$ 는 ``$132$ 로는 못 이겼고 $3$ 으로 줄이면 다시 못 이겼다.
최적점은 다섯''으로 적었다. 지금 자료에서

\begin{center}
$5$ 개 $\Delta=-0.0579$ \qquad $3$ 개 $\Delta=-0.0568$
\end{center}

\textbf{사실상 같고 둘 다 본페로니 $99\%$ 를 통과한다.} 봉우리라고 부를
근거가 없다.

\begin{quote}
\textbf{$U$ 자의 양끝이 아니라 바닥이 넓다.} 확실한 것은 ``\textbf{피처를
많이 주면 못 이기고 열 개 아래로 줄이면 이긴다}''이지 ``다섯이 최적''이
아니다. $\Delta$ 차이 $0.0011$ 은 구간 폭 $0.08$ 의 $1/70$ 이다 ---
\textbf{봉우리를 읽은 것은 잡음을 읽은 것}이다(노트 $608$ 이 \texttt{lab/}
쪽 도메인 표에서 겪은 것과 같은 자리).
\end{quote}

\section{두 갈래가 같은 다섯으로 수렴했다}

논문 $4$ 가 남긴 다섯은 타깃 폭 $\cdot$ 매장 노출도 $\cdot$ 입장 허들
$\cdot$ 미디어 투입 $\cdot$ 굿즈 규모다. \textbf{\texttt{lab/} 의
\texttt{ALL5} 와 정확히 같은 다섯}이다 --- \texttt{target\_breadth} $\cdot$
\texttt{venue\_prominence} $\cdot$ \texttt{entry\_friction} $\cdot$
\texttt{media\_push} $\cdot$ \texttt{goods\_scale}.

\begin{center}
\begin{tabular}{lrr}
\hline
 & \texttt{lab/} & 팝업 갈래 \\
\hline
표본 & $3{,}369$ 행 $\cdot$ $11$ 도메인 & $75$ 행 $\cdot$ 계수 계열 \\
자 & 순위 $\rho$ & MAE(로그 일평균) \\
결론 & 안 본 시장에 전이 $82\!\sim\!95\%$ & 상수를 $0.058$ 이긴다 \\
\hline
\end{tabular}
\end{center}

\textbf{두 갈래는 자료도 자도 다르다.} 같은 다섯이 양쪽에서 남은 것은
독립적인 두 증거다 --- 다만 \textbf{같은 사람이 같은 태거로 매긴 축}이므로
완전히 독립은 아니다(축 정의가 공유된다).

\section*{조항}

\begin{quote}
\textbf{내 의심이 틀리는 것도 결과다.} 군집과 선택 둘 다 원리상 옳은
의심이었고 둘 다 이 자료에서는 안 물었다. \textbf{그 대신 안 의심하던 것
(``최적점은 다섯'')이 무너졌다.} 자를 겨눌 때 어디가 무너질지는 미리
모른다 --- \textbf{그래서 겨눠 보는 것이다.}
\end{quote}

\section*{남은 것}

\begin{itemize}
\item \textbf{채점 행이 $40$ 개다.} $n{=}75$ 중 $35$ 는 rolling-origin
의 최초 학습 구간이라 채점에 안 들어간다. \textbf{$\Delta$ 의 구간 폭이
$0.08$ 인 이유가 여기 있다.}
\item \textbf{\texttt{paper/figs.py} 의 수치가 하드코딩이다.} 그 파일
머리말은 ``데이터에서 직접 계산하므로 라벨이 바뀌면 그림도 바뀐다''인데
\texttt{fig\_feature\_curve} 는 리터럴 다섯 줄이다. \textbf{라벨이 바뀌어도
그림이 안 바뀐다} --- 이번에 실제로 $-0.0540$ 이 $-0.0579$ 로 움직였다.
\item \textbf{옛 헤드라인은 재현 안 된다} --- 커밋 \texttt{e24918f} 의
``내부 레코드 $n{=}73$, $0.378\!\to\!0.329$''는 지금 자료에서
$0.375\!\to\!0.356$ 이고 구간이 $0$ 을 문다. 그 뒤 라벨 재채굴이 여러 번
있었으므로 \textbf{자료가 바뀐 것}이지만, \textbf{옛 수치를 인용하면 안
된다.}
\end{itemize}

\end{document}
